\documentclass{article}


\usepackage{arxiv}

\usepackage[utf8]{inputenc} % allow utf-8 input
\usepackage[T1]{fontenc}    % use 8-bit T1 fonts
\usepackage{hyperref}       % hyperlinks
\usepackage{url}            % simple URL typesetting
\usepackage{booktabs}       % professional-quality tables
\usepackage{amsfonts}       % blackboard math symbols
\usepackage{nicefrac}       % compact symbols for 1/2, etc.
\usepackage{microtype}      % microtypography
\usepackage{lipsum}
\usepackage{graphicx}
\usepackage[brazil]{babel}
\graphicspath{ {./images/} }


\title{Modelagem e controle de um conversor flyback por meio de estratégias de controle preditivo.}


\author{
 Ricardo Augusto de Araújo Machado \\
  Programa de Pós-Graduação em Engenharia Elétrica e de Computação\\
  Universidade Federal da Bahia\\
  \texttt{machado.ricardo@ufba.br} \\
  %% examples of more authors
  %% \AND
  %% Coauthor \\
  %% Affiliation \\
  %% Address \\
  %% \texttt{email} \\
  %% \And
  %% Coauthor \\
  %% Affiliation \\
  %% Address \\
  %% \texttt{email} \\
  %% \And
  %% Coauthor \\
  %% Affiliation \\
  %% Address \\
  %% \texttt{email} \\
}

\begin{document}
\maketitle
\begin{abstract}
  Escrever resumo.
\end{abstract}


% keywords can be removed
%\keywords{First keyword \and Second keyword \and More}


\section{Introdução}
O conversor flyback é um conversor estático de corrente contínua que opera tanto como elevador quanto como abaixador de tensão. Por conta da isolação galvânica presente no dispostivo, essa topologia de circuito é extremamente popular em fontes chaveadas(SMPS) de baixa potência \cite{Picard2010}.

A operação do conversor flyback consiste em armazenar energia na indutância de magnetização do transformador enquanto o transistor conduz corrente e, em seguida, transferir essa energia armazenada para a saída enquanto o transistor estiver aberto \cite{Picard2010}.

Mesmo com a popularidade desta topologia, o controle do conversor flyback apresenta uma determinada complexidade, pois o sistema é não linear, chaveado e apresenta um zero de fase não mínima. Considerando os desafios citados, o controle preditivo baseado em modelo(MPC) é uma metodologia de controle adequada para lidar com as complexidades do conversor.


\section{Modelagem e dimensionamento do conversor flyback.}

\subsection{Especificações escolhidas para o conversor flyback.}

O conversor flyback é dimensionado para operar no modo de condução contínua, ou seja, a corrente no indutor é projetada para ser sempre maior que zero durante o ciclo de chaveamento. Ademais, a frequência de chaveamento é definida como \(50\ \mathrm{kHz}\) e a potência de saída nominal é \(6\ W\). As especificações completas são listadas na Tabela




\subsection{Modelagem dinâmica do conversor.}

A modelagem do conversor flyback, descrita nos livros \cite{Erickson2020} e, consiste na determinação das equações dinâmicas da tensão no indutor e da corrente no capacitor durante os dois estados de operação do conversor, chave fechada e chave aberta.

\subsection{Dimensionamento do conversor.}

As expressões para a tensão média no capacitor e para a corrente média no indutor, dadas por (\ref{eq:tensao_media}) e (\ref{eq:corrente_media}), são obtidas por meio do balanço Volt-Segundo no indutor e do balanço de carga no capacitor, respectivamente.

\begin{equation}
  V_\mathrm{C}
  \label{eq:tensao_media}
\end{equation}

\begin{equation}
  I_\mathrm{Lm}
  \label{eq:corrente_media}
\end{equation}

A razão cíclica de equilibrio é determinada por meio de . Para as tensões nominais definidas, a razão cícilica obtida é \(\)



\bibliographystyle{unsrt}
\bibliography{references}




\end{document}
